\section{Fine-Tuning with Cut-and-Paste Outlier Exposure}
\label{sec:finetuning}

In this work, we fine-tune the EoMT segmentation model using a
\textit{cut-and-paste outlier exposure} strategy. The goal is to expose the model
to synthetic object-level anomalies during training, while preserving its
semantic segmentation performance on in-distribution Cityscapes data.

The method constructs augmented training samples by pasting transparent object
cutouts extracted from COCO into Cityscapes images. The pasted regions are treated
as out-of-distribution pixels, while the remaining pixels retain their original
Cityscapes semantic labels. This allows the model to continue learning the
in-distribution segmentation task while also being regularized to produce
uncertain predictions on anomalous regions.

\subsection{Outlier Injection Strategy}

\paragraph{COCO Object Cutouts.}
Outlier objects are extracted from COCO instance annotations as RGBA cutouts. We
use a manually selected subset of COCO categories that are plausible as road-scene
anomalies, including animals, bags, bottles, furniture, plants, and other
portable objects. For each selected instance, the COCO segmentation mask is used
to preserve only the foreground object and store it with an alpha channel.

During training, one cutout is sampled from the generated outlier pool and pasted
into a Cityscapes image. The current implementation samples from the available
cutout files rather than using an adaptive or learned category weighting scheme.

\paragraph{Scale Prior.}
To expose the model to anomalies of different sizes, the pasted object area is
sampled from a bimodal distribution. In the reported configuration, small
outliers are sampled with high probability:
\begin{equation}
P(s) =
\begin{cases}
\mathcal{U}(0.0005, 0.008) & \text{with probability } 0.85, \\
\mathcal{U}(0.02, 0.08) & \text{with probability } 0.15.
\end{cases}
\end{equation}

Here, $s$ denotes the target area ratio relative to the full image. This choice
biases training toward small road obstacles while still occasionally introducing
larger object-level anomalies.

\paragraph{Perspective-Aware Placement.}
The sampled scale is modulated according to the vertical placement of the object
in the image:
\begin{equation}
s_{\text{eff}} =
s \cdot \max\left(0.1, \left(\frac{y}{H}\right)^{\gamma}\right),
\end{equation}
where $y$ is the vertical coordinate of the paste center, $H$ is the image
height, and $\gamma$ controls the perspective strength. This makes objects pasted
near the top of the image smaller than objects pasted in the foreground.

In the current configuration, placement is primarily road-biased. With high
probability, the paste center is sampled from pixels belonging to the Cityscapes
road class, preferably in the lower part of the image. If no suitable road pixel
is selected, the object is placed at a random location in the lower image region.

The implementation also supports boundary-aware placement near selected semantic
classes such as buildings, walls, fences, poles, vegetation, and terrain.
However, this option is disabled in the reported configuration.

\paragraph{Photometric Blending.}
Before pasting, the cutout is resized according to the sampled effective scale.
A color jitter transformation is applied to increase appearance diversity. The
implementation also supports edge feathering to reduce sharp mask artifacts at
object boundaries. In the reported configuration, hard alpha blending is used
with feathering enabled and no additional global noise.

\subsection{Optimization Objective}

Training uses a combined objective with a semantic segmentation term and an OOD
regularization term:
\begin{equation}
\mathcal{L} =
\lambda_{CE} \mathcal{L}_{CE} +
\lambda_{OOD} \mathcal{L}_{OOD}.
\end{equation}

The cross-entropy loss $\mathcal{L}_{CE}$ is computed only on valid
in-distribution Cityscapes pixels. Pasted OOD pixels are assigned the ignore
label for the semantic cross-entropy term, so they do not contribute as ordinary
Cityscapes classes.

The OOD loss is computed only over pasted anomaly pixels. In the reported
configuration, $\mathcal{L}_{OOD}$ combines a maximum-entropy objective with a
logit-norm penalty:
\begin{equation}
\mathcal{L}_{OOD}
=
\lambda_H \mathcal{L}_{H}
+
\lambda_N \mathcal{L}_{N}.
\end{equation}

The entropy term encourages high softmax uncertainty on pasted outliers, while
the logit-norm term discourages overconfident high-magnitude logits in OOD
regions. The OOD loss weight is gradually warmed up during the first training
epochs to reduce instability.

During fine-tuning, the DINOv2 encoder of EoMT is frozen. Only the remaining
segmentation components are updated. This design is intended to preserve the
pretrained visual representation and reduce catastrophic forgetting, while still
adapting the segmentation head to recognize synthetic outlier regions as
uncertain.

\subsection{Empirical Behavior of the Fine-Tuned Model}

The fine-tuned model is evaluated on multiple road-anomaly benchmarks using
standard OOD detection metrics such as AuPRC and FPR@TPR95. The expected effect
of the fine-tuning procedure is to improve the separation between
in-distribution road-scene pixels and object-like anomalous regions.

Because the training anomalies are structured object-level cutouts, the method is
most directly aligned with medium and large object-like anomalies. Improvements
are therefore expected to be more visible when the test anomalies are spatially
coherent and visually similar to pasted objects. Gains may be more limited for
very small, ambiguous, or texture-like anomalies.

The frozen encoder and retained cross-entropy supervision help preserve
in-distribution semantic segmentation behavior while the OOD regularization
encourages less confident predictions on pasted anomaly pixels.

\subsection{Summary}

The proposed fine-tuning strategy introduces three main mechanisms:

\begin{itemize}
    \item synthetic object-level outlier exposure using COCO instance cutouts,
    \item perspective-aware and road-biased placement of pasted anomalies,
    \item OOD regularization through entropy maximization and logit-norm control.
\end{itemize}

Overall, this approach adapts EoMT toward improved OOD awareness without changing
the underlying architecture. It uses synthetic anomalies only during training and
preserves the standard semantic segmentation output at inference time.